\documentclass[10pt,a4paper]{amsart}
\usepackage[margin=19mm]{geometry}
\usepackage{amsmath,amssymb,mathtools}
\usepackage[scr=rsfs,cal=euler]{mathalfa}
\usepackage{xfrac}
\usepackage[unicode,hidelinks,bookmarks=false]{hyperref}
\newcommand{\ie}{i.e.\ }
\newcommand{\cf}{cf.\ }
\newcommand{\resp}{respectively\ }
\newcommand{\Subsection}{\subsection}
\providecommand{\nobreakdash}{\nobreak\hbox{-}}
\setlength{\emergencystretch}{2em}
\multlinegap0pt
\usepackage{bm}
\usepackage{bbm}
\usepackage{dsfont}
\usepackage{xspace}
\usepackage{cancel}
\usepackage{mathtools}
\usepackage{amsthm}
\makeatletter
\@ifundefined{theorem}{\newtheorem{theorem}{Theorem}}{}
\@ifundefined{lemma}{\newtheorem{lemma}[theorem]{Lemma}}{}
\@ifundefined{proposition}{\newtheorem{proposition}[theorem]{Proposition}}{}
\makeatother
\newcommand{\Nrm}[2]{\mathbf{N}_{#1}(#2)}
\newcommand{\Q}{\mathbb{Q}}
\newcommand{\Z}{\mathbb{Z}}
\title[The distribution of $k$-free ideals in ray class groups]{The distribution of $k$-free ideals in ray class groups}
\author{Adrian Barquero-Sanchez, Jack Heimrath, Bernd Sing, Nicol\'as Sirolli, Caylee Spivey, Michael Wijaya}
\date{Source: arXiv:2608.13509v1 (CC BY 4.0)}
\begin{document}
\maketitle

\noindent\emph{Source and licence.} The distribution of $k$-free ideals in ray class groups. Version arXiv:2608.13509v1 (CC BY 4.0), distributed under the Creative Commons Attribution 4.0 International licence, as recorded in the arXiv licence metadata for this exact version and shown on the abstract page https://arxiv.org/abs/2608.13509v1. Full rights notice: licenses/arxiv-2608.13509-CC-BY-4.0.txt.\par\smallskip

\noindent\textit{Editorial scope.} Counted unit: the Theorem whose source label is \texttt{thm:main-theorem-tuples} in the article (source0.tex lines 1969--2022, source label \texttt{thm:main-theorem-tuples}), delivered together with the article's own complete original proof of it (source0.tex lines 2024--2214, one \texttt{proof} environment of 191 lines). No mathematics is written, repaired or re-derived: statement and proof are the article's own text line for line. Required context retained uncounted: prop:Q-identity-with-moebius (lines 1571--1587); lem:several-cases-estimation (lines 1733--1745); lem:finite-sum-with-moebius-to-dedekind-zeta (lines 1863--1878); eqn:unified-ideal-counting (lines 1955--1964). Every definition, display and label of those lines is retained unchanged. Omitted from this package and not redistributed: 1--1570, 1588--1732, 1746--1862, 1879--1954, 1965--1968, 2023--2023, 2215--2432. Nothing inside a retained excerpt is omitted or altered: every retained body is delivered exactly as the article's lines are written, so no omission receipt of any kind accompanies this package. Citations: the retained excerpts carry no citation command at all, so no bibliography entry of the article is transcribed and no reference is redistributed. Reference adaptation: none was needed and none was made; every reference inside the delivered unit resolves inside it, so no unresolved reference or citation is produced. Printed slips kept literal: no printed word, symbol, hypothesis or number of the counted statement or of its proof was altered. Layout and class: the article's own class is \texttt{amsart} with packages including amscd, babel, inputenc, amsmath, amssymb, enumitem, booktabs, comment, xcolor, hyperref, tabu, siunitx, orcidlink, pdflscape; the wrapper is the lane's pinned \texttt{amsart} editorial wrapper at 10pt on A4, which restates the theorem-like environments the retained text needs and adds the article's own macro definitions that the retained text needs (\texttt{\textbackslash Nrm}, \texttt{\textbackslash Q}, \texttt{\textbackslash Z}), with extra wrapper packages (bm, bbm, dsfont, xspace, cancel, mathtools). The delivered numbering is the unit's own. Assets: no retained span contains a figure environment, a graphics-inclusion command, a PDF-page inclusion, a table or a TikZ picture; no image file is copied into this package, the package declares no asset dependency and no image file is redistributed with it. Licence: version arXiv:2608.13509v1 of this article is distributed under the Creative Commons Attribution 4.0 International licence, as recorded in the arXiv licence metadata for this exact version and shown on the abstract page https://arxiv.org/abs/2608.13509v1. A full rights notice is delivered at licenses/arxiv-2608.13509-CC-BY-4.0.txt. Characters: every retained excerpt is pure ASCII, so the delivered engine, XeLaTeX with its default Latin Modern text font, reports no missing character.
\begin{proposition}\label{prop:Q-identity-with-moebius}
Let $k,r\in\Z_{\geq 1}$ and let $\boldsymbol{\mathfrak{A}} = (\mathfrak{A}_1,\dots,\mathfrak{A}_r) \in \left(\mathrm{Cl}_K^{\mathfrak{m}}\right)^r$.
Then, for every $X>0$, we have
\begin{align*}
Q_K^{k,r}(X;\boldsymbol{\mathfrak{A}})
&=
\sum_{\mathfrak{a}\in
\mathcal{I}_K^{\mathfrak{m}}(\sqrt[k]{X})}
\mu_K(\mathfrak{a})
\prod_{i=1}^{r}
\mathcal{H}_K^{\mathfrak{m}}
\left(
\frac{X}{\Nrm{K/\Q}{\mathfrak{a}^k}};
\mathfrak{A}_i\cdot[\mathfrak{a}^{-k}]
\right).
\end{align*}
\end{proposition}

\begin{lemma}\label{lem:several-cases-estimation}
Let $s\geq 0$ and $T > 1$. Then
\begin{align*}
\sum_{\mathfrak{a}\in\mathcal{I}_K^{\mathfrak{m}}(T)}
\frac{1}{\Nrm{K/\Q}{\mathfrak{a}}^s}
=
\begin{cases}
O(T^{1-s}), & \text{if $0\leq s<1$},\\[4pt]
O(\log T), & \text{if $s=1$},\\[4pt]
O(1), & \text{if $s>1$}.
\end{cases}
\end{align*}
\end{lemma}

\begin{lemma}\label{lem:finite-sum-with-moebius-to-dedekind-zeta}
Let $s>1$ and $T\geq 1$. Then
\begin{align*}
\sum_{\mathfrak{a}\in\mathcal{I}_K^{\mathfrak{m}}(T)}
\frac{\mu_K(\mathfrak{a})}
{\Nrm{K/\Q}{\mathfrak{a}}^s}
&=
\frac{1}{\zeta_K(s)}
\prod_{\mathfrak{p}\mid\mathfrak{m}_0}
\left(
1-\frac{1}{\Nrm{K/\Q}{\mathfrak{p}}^s}
\right)^{-1}
+
O(T^{1-s}).
\end{align*}
\end{lemma}

\begin{equation}
\label{eqn:unified-ideal-counting}
\mathcal{H}_K^{\mathfrak{m}}(X;\mathfrak{A})
=
\rho_K^{\mathfrak{m}}X
+
O\left(
X^{1-\delta_d(\mathfrak{m})}
\right).
\end{equation}

\begin{theorem}\label{thm:main-theorem-tuples}
Let $K$ be a number field of degree $[K:\Q]=d$, let
$\mathfrak{m}$ be a modulus for $K$, and let $k,r\in\Z_{\geq1}$
be such that $(k,r)\neq(1,1)$. Let
\[
\boldsymbol{\mathfrak{A}}
=
(\mathfrak{A}_1,\ldots,\mathfrak{A}_r)
\in
\left(\mathrm{Cl}_K^{\mathfrak{m}}\right)^r.
\]
Then
\begin{align}
Q_K^{k,r}(X;\boldsymbol{\mathfrak{A}})
&=
\frac{\left(\rho_K^{\mathfrak{m}}\right)^r}
{\zeta_K(rk)}
\prod_{\mathfrak{p}\mid\mathfrak{m}_0}
\left(
1-\frac{1}
{\Nrm{K/\Q}{\mathfrak{p}}^{rk}}
\right)^{-1}
X^r
+
E_{K,\mathfrak{m},k,r}(X),
\label{eqn:main-theorem-tuples-asymptotic}
\end{align}
where
\begin{equation}\label{eqn:main-theorem-tuples-error}
E_{K,\mathfrak{m},k,r}(X)
=
\begin{cases}
\displaystyle
O\left(X^{1/k}\right),
&
k\left(r-\delta_d(\mathfrak{m})\right)<1,
\\[8pt]
\displaystyle
O\left(
X^{r-\delta_d(\mathfrak{m})}\log X
\right),
&
k\left(r-\delta_d(\mathfrak{m})\right)=1,
\\[8pt]
\displaystyle
O\left(
X^{r-\delta_d(\mathfrak{m})}
\right),
&
k\left(r-\delta_d(\mathfrak{m})\right)>1.
\end{cases}
\end{equation}

\end{theorem}

\begin{proof}
For simplicity, write
\[
\delta:=\delta_d(\mathfrak{m})
\qquad\text{and}\qquad
\mathcal{J}
=
\mathcal{I}_K^{\mathfrak{m}}\left(X^{1/k}\right).
\]
By Proposition~\ref{prop:Q-identity-with-moebius}, we have
\begin{align}\label{eqn:Q-equality-with-moebius-for-proof}
Q_K^{k,r}(X;\boldsymbol{\mathfrak{A}})
&=
\sum_{\mathfrak{a}\in\mathcal{J}}
\mu_K(\mathfrak{a})
\prod_{i=1}^{r}
\mathcal{H}_K^{\mathfrak{m}}
\left(
\frac{X}{\Nrm{K/\Q}{\mathfrak{a}^k}};
\mathfrak{A}_i\cdot[\mathfrak{a}^{-k}]
\right).
\end{align}

Since the ray class group $\mathrm{Cl}_K^{\mathfrak{m}}$ is finite,
the implied constant in \eqref{eqn:unified-ideal-counting} may be
chosen uniformly over all ray classes. Hence, for any
$\mathfrak{C}_1,\ldots,\mathfrak{C}_r
\in\mathrm{Cl}_K^{\mathfrak{m}}$ and any $Y\geq1$,
\begin{align}
\prod_{i=1}^{r}
\mathcal{H}_K^{\mathfrak{m}}(Y;\mathfrak{C}_i)
&=
\prod_{i=1}^{r}
\left(
\rho_K^{\mathfrak{m}}Y
+
O\left(Y^{1-\delta}\right)
\right)
\notag\\
&=
\left(\rho_K^{\mathfrak{m}}\right)^rY^r
+
O\left(Y^{r-\delta}\right).
\label{eqn:product-H-asymptotic-unified}
\end{align}
Indeed, every term other than the main product has the form
\[
\left(\rho_K^{\mathfrak{m}}Y\right)^t
O\left(Y^{1-\delta}\right)^{r-t}
\]
for some $0\leq t\leq r-1$, and hence has order
\[
O\left(
Y^{t+(r-t)(1-\delta)}
\right)
=
O\left(
Y^{r-(r-t)\delta}
\right).
\]
The largest exponent occurs for $t=r-1$, giving
$O(Y^{r-\delta})$.

Next, applying \eqref{eqn:product-H-asymptotic-unified} with the choices
\[
Y
=
\frac{X}{\Nrm{K/\Q}{\mathfrak{a}^k}} \quad \text{and} \quad \mathfrak{C}_i
=
\mathfrak{A}_i\cdot[\mathfrak{a}^{-k}]
\]
we obtain
\begin{align}\label{eqn:product-H-asymptotic-substitution}
\prod_{i=1}^{r}
\mathcal{H}_K^{\mathfrak{m}}
\left(
\frac{X}{\Nrm{K/\Q}{\mathfrak{a}^k}};
\mathfrak{A}_i\cdot[\mathfrak{a}^{-k}]
\right) =
\left(\rho_K^{\mathfrak{m}}\right)^r
\frac{X^r}
{\Nrm{K/\Q}{\mathfrak{a}}^{rk}}
+
O\left(
\frac{X^{r-\delta}}
{\Nrm{K/\Q}{\mathfrak{a}}^{k(r-\delta)}}
\right).
\end{align}
Then, using \eqref{eqn:product-H-asymptotic-substitution} in
\eqref{eqn:Q-equality-with-moebius-for-proof} yields
\begin{align}
Q_K^{k,r}(X;\boldsymbol{\mathfrak{A}})
&=
\left(\rho_K^{\mathfrak{m}}\right)^r
\left( \sum_{\mathfrak{a}\in\mathcal{J}}
\frac{\mu_K(\mathfrak{a})}
{\Nrm{K/\Q}{\mathfrak{a}}^{rk}}
\right)X^r + O\left( \sum_{\mathfrak{a}\in\mathcal{J}}
\frac{1}{\Nrm{K/\Q}{\mathfrak{a}}^{k(r-\delta)}}\right)
X^{r-\delta}.
\label{eqn:Q-expression-with-two-sums}
\end{align}

We first consider the sum containing the M\"obius function.
Since $(k,r)\neq(1,1)$, we have $rk>1$. Therefore, using
Lemma \ref{lem:finite-sum-with-moebius-to-dedekind-zeta} with the choices $s=rk$ and $T=X^{1/k}$ gives
\begin{align}
\sum_{\mathfrak{a}\in\mathcal{J}}
\frac{\mu_K(\mathfrak{a})}
{\Nrm{K/\Q}{\mathfrak{a}}^{rk}}
&=
\frac{1}{\zeta_K(rk)}
\prod_{\mathfrak{p}\mid\mathfrak{m}_0}
\left(
1-\frac{1}
{\Nrm{K/\Q}{\mathfrak{p}}^{rk}}
\right)^{-1}
+
O\left(X^{1/k - r}\right).
\end{align}
Consequently,
\begin{align}
\left(\rho_K^{\mathfrak{m}}\right)^r
\left(
\sum_{\mathfrak{a}\in\mathcal{J}}
\frac{\mu_K(\mathfrak{a})}
{\Nrm{K/\Q}{\mathfrak{a}}^{rk}}
\right)X^r =
\frac{\left(\rho_K^{\mathfrak{m}}\right)^r}
{\zeta_K(rk)}
\prod_{\mathfrak{p}\mid\mathfrak{m}_0}
\left(
1-\frac{1}
{\Nrm{K/\Q}{\mathfrak{p}}^{rk}}
\right)^{-1} X^r + O\left(X^{1/k}\right).
\label{eqn:main-sum-asymptotic}
\end{align}

We next estimate the second sum in \eqref{eqn:Q-expression-with-two-sums}. Set $s:=k(r-\delta)$. Since $r\geq1$ and $0<\delta\leq1$, we have $s\geq0$, so
Lemma \ref{lem:several-cases-estimation} applies with
$T=X^{1/k}$. Hence
\begin{align}
&X^{r-\delta}
\sum_{\mathfrak{a}\in\mathcal{J}}
\frac{1}
{\Nrm{K/\Q}{\mathfrak{a}}^{k(r-\delta)}}
=
\begin{cases}
\displaystyle
O\left(
X^{r-\delta}
X^{\frac{1-k(r-\delta)}{k}}
\right)
=
O\left(X^{1/k}\right),
&
k(r-\delta)<1,
\\[10pt]
\displaystyle
O\left(
X^{r-\delta}\log X
\right),
&
k(r-\delta)=1,
\\[10pt]
\displaystyle
O\left(
X^{r-\delta}
\right),
&
k(r-\delta)>1.
\end{cases}
\label{eqn:secondary-error-cases}
\end{align}


It remains only to compare these estimates with the error
$O(X^{1/k})$ arising from the main sum in
\eqref{eqn:main-sum-asymptotic}. We claim that, in each of the three
cases above, the error term coming from the second sum in
\eqref{eqn:Q-expression-with-two-sums} is at least as large as
$O(X^{1/k})$, and therefore determines the final error term. Indeed,
if $k(r-\delta)<1$, the two error terms are both $O(X^{1/k})$; if
$k(r-\delta)=1$, then $r-\delta=1/k$, so
$O(X^{1/k})$ is absorbed by
$O(X^{r-\delta}\log X)$; and if $k(r-\delta)>1$, then
$r-\delta>1/k$, so $O(X^{1/k})$ is absorbed by
$O(X^{r-\delta})$.

Combining these estimates proves \eqref{eqn:main-theorem-tuples-asymptotic} and \eqref{eqn:main-theorem-tuples-error}.
\end{proof}

\end{document}
